\documentclass[a4paper]{article}
\usepackage{geometry}
\usepackage{booktabs,tabularx}
\usepackage[round,authoryear]{natbib}
\usepackage[hidelinks]{hyperref}

\title{Recoverability of causal structure\\from psychological networks?}
\author{Harsh Nagouda \\ \small BEP research proposal, first draft}
\date{21 September 2026}

\begin{document}
\maketitle

\section{Introduction}

Psychological network models represent symptoms as nodes and their associations as
edges, while the theories that motivate them concern something different: the processes
by which one symptom brings about another. The gap between the two rarely matters until
a cross-sectional network is read as a map of where to intervene, at which point an
association is being asked to carry a causal claim it was never estimated to support.
This project asks how much causal structure the cross-sectional datasets collected in
ReBayesed can actually determine, and what, in the data or in the methods, sets the
limit on orienting an edge.

\paragraph{Research problem.} In these networks an edge is a partial correlation: an
association that survives conditioning on every other measured symptom
\citep{borsboom2013}. Causal discovery narrows the possible orientations of such edges
using the pattern of conditional independencies, but observational data typically leave
several structures compatible with the same pattern. How much ambiguity remains has not
been assessed systematically across a published corpus of psychological networks.

\paragraph{Relevance.} Undirected networks are routinely used to discuss intervention,
yet intervention requires direction. The ReBayesed export shows how thin the evidence
often is even for the undirected edges: in the tenth of networks with the least data per
parameter, 56.7\% of edges have inconclusive Bayesian evidence, against 17.1\% in the
tenth with the most. If an edge cannot be established as present, its direction is
unlikely to fare better, and I will test whether that expectation holds.

\paragraph{Significance.} Of the 294 networks, 147 have fewer than 5 observations per
edge parameter and 44 have fewer than 1. If orientation tracks data volume, the result
gives sample-planning guidance for future studies; if orientation stays ambiguous even
where edge evidence is strong, the result identifies a limit that more data would not
remove. Either outcome is informative.

\paragraph{Novelty.} \citet{huth2026} quantified the evidence for undirected edges;
\citet{huber2024} set out why causal structure is often not identifiable from
observational data. Neither applied the other's question to this corpus. This project
does.

\section{Related work}

\citet{huth2026} re-analysed 294 cross-sectional psychological networks from 126 papers
with Bayesian model averaging and showed that edge evidence depends strongly on sample
size. Their export is the input here, because it lets me ask of the same datasets
whether the evidence that supports an edge also supports a direction.

\citet{huber2024} explains why observational data can leave causal structures
indistinguishable, particularly when common causes are unmeasured. I use this framework
to fix the assumptions under which any orientation the analysis produces should be read.

For the analysis I will use \texttt{causal-learn} \citep{zheng2024}, which implements
both PC and FCI in Python. PC assumes that all common causes are measured; FCI drops
that assumption at the cost of weaker conclusions. Running both from one library keeps
the comparison in RQ3 consistent, since any difference then reflects the assumption
rather than the implementation.

\section{Research questions}

\textbf{Main question: how much causal structure can be recovered from these psychological
datasets, and what limits that recovery?}

\begin{enumerate}
\item Do networks with more observations per edge parameter yield a greater proportion of determined edge
orientations?
\item Does stronger Bayesian evidence for an edge translate into a determined direction
for that edge?
\item How do the conclusions change when unmeasured common causes are allowed?
\end{enumerate}

As a supplementary comparison, I will check whether independent studies that measured
the same variables agree on orientation.

\section{Expected outcomes}

\paragraph{Approach.} The export contains correlation matrices, sample sizes and
metadata rather than participant-level observations, so the first task is to confirm
that Fisher's z tests, and \texttt{causal-learn} on top of them, can run from a
correlation matrix and a sample size alone. The export also needs cleaning before
anything runs on it: 233 of the 294 networks have variable names that disagree with
their matrix labels, and the export holds 294 networks where the article reports 293.
The data are 71 continuous, 196 ordinal and 27 mixed networks. Exploratory analysis of
the export is complete and is the source of these counts; the cleaning itself has not
started. A pilot in week 5 will decide which continuous networks form the analysis
subset and which test settings to fix.

\paragraph{RQ1: data volume and orientation.} I will run PC on each network in the
subset, record its adjacencies and the proportion of edges it orients, and then check
how sensitive those proportions are to the significance threshold and to variable
ordering. Because 58 of the 126 papers report several networks that are potentially
dependent, I will select one network per paper by a prespecified rule and report the
others separately, so that dependent networks do not inflate the trend. I expect the
oriented proportion to rise with observations per edge parameter but to stay well short
of complete, which is the pattern that would show data volume matters without being the
whole story.

\paragraph{RQ2: statistical evidence and direction.} I will relate each network's
oriented proportion to its proportion of inconclusive Bayesian edges under each of the
three priors. Bayesian edge evidence enters as a covariate rather than as ground truth,
because a partial-correlation edge and a causal adjacency are not the same object: an
edge can be estimated with confidence and still leave its direction undetermined. Those
edges, strong evidence with no direction, are the ones I will examine most closely,
since they mark where the limit is set by identifiability rather than by sample size.

\paragraph{RQ3: allowing unmeasured common causes.} I will run FCI on the same subset.
FCI and PC do not output the same objects, so counting arrowheads would compare unlike
things; instead I will compare statements both can make, such as whether one variable
is an ancestor of another, and quantify how much of PC's structure dissolves once
unmeasured common causes are permitted. Both algorithms still assume acyclicity, the
Markov and faithfulness conditions, and the absence of selection bias
\citep{huber2024}, and these assumptions bound what any result in this project can
claim.

\paragraph{Cross-study comparison and outputs.} Depression is the largest subtopic, with
32 networks, and the first candidate for comparing orientations across studies; whether
that comparison is feasible depends on the metadata cleaning in week 4. Where studies
disagree, I will distinguish a direction left unresolved from one that conflicts across
studies, since only the second bears on replication. The deliverables are a cleaned
pipeline, reproducible analyses, figures and a written interpretation that claims no
more than the data determine under the assumptions above.

\section{Research plan}

The project runs for 13 weeks and is now in week 4. The first three weeks went to
exploratory analysis of the export, which produced the descriptive figures above and
exposed the cleaning problems; the plan below covers the remaining ten weeks. Data
coordination is shared with the other students on the project; I lead the Python
analyses and the interpretation. Literature review and writing run alongside the
analysis rather than after it.

\begin{center}
\begin{tabularx}{\textwidth}{@{}r X@{}}
\toprule
\textbf{Week} & \textbf{Activity and output} \\
\midrule
4 & Clean metadata; resolve the 294--293 discrepancy; identify eligible networks and studies. \\
5 & Pilot: validate Fisher's z and \texttt{causal-learn} on correlation-matrix input; fix the subset and test settings. \\
6 & Build and validate the PC pipeline; run PC on the subset. \\
7 & Summarise adjacencies and orientations; analyse data volume and sensitivity (RQ1). \\
8 & Relate orientation to Bayesian evidence (RQ2). \\
9 & Run FCI on the matched subset. \\
10 & Compare analyses with and without unmeasured causes (RQ3). \\
11 & Complete the cross-study comparison and checks; draft results. \\
12 & Complete the thesis draft and discuss feedback. \\
13 & Revise, check reproducibility and prepare the presentation. \\
\bottomrule
\end{tabularx}
\end{center}

The compression falls on weeks 6 and 11, where building the PC pipeline shares a week
with running it and the cross-study comparison shares a week with the final checks.
Both are workable if the pilot in week 5 settles the subset and settings, which is why
that week is protected. Two further constraints are addressed early. If FCI proves
computationally prohibitive, I will run both algorithms on a smaller prespecified
subset so that RQ3 survives intact. If no independent studies share variables, I will
report that as a limitation of the corpus rather than substitute within-paper
comparisons, which cannot count as replication. The subset is fixed by week 5.

\bibliographystyle{plainnat}
\bibliography{references}

\end{document}